\documentclass[10pt,a4paper]{amsart}
\usepackage[margin=19mm]{geometry}
\usepackage{amsmath,amssymb,mathtools}
\usepackage[scr=rsfs,cal=euler]{mathalfa}
\usepackage{xfrac}
\usepackage[unicode,hidelinks,bookmarks=false]{hyperref}
\newcommand{\ie}{i.e.\ }
\newcommand{\cf}{cf.\ }
\newcommand{\resp}{respectively\ }
\newcommand{\Subsection}{\subsection}
\providecommand{\nobreakdash}{\nobreak\hbox{-}}
\setlength{\emergencystretch}{2em}
\multlinegap0pt
\usepackage{bm}
\usepackage{algorithm}\usepackage{algpseudocode}
\usepackage{tcolorbox}
\newtheorem{proposition}{Proposition}
\title{Hierarchical Mixture-of-Experts with Two-Stage Optimization}
\author{Gleb Molodtsov, Alexander Miasnikov, Aleksandr Beznosikov}
\date{Source: arXiv:2605.08292v1 (CC BY 4.0)}
\begin{document}
\maketitle

\noindent\emph{Source and licence.} Hierarchical Mixture-of-Experts with Two-Stage Optimization. Version arXiv:2605.08292v1 (CC BY 4.0), distributed by arXiv under the Creative Commons Attribution 4.0 International licence, http://creativecommons.org/licenses/by/4.0/, as recorded in the arXiv licence metadata for this exact version and shown on the abstract page https://arxiv.org/abs/2605.08292v1. Full licence notice: \texttt{licenses/arxiv-2605.08292-CC-BY-4.0.txt}.\par\smallskip

\noindent\textit{Editorial scope.} Counted unit: the article's proposition thm:cv-l2 (source0.tex lines 1137--1144). The article's complete original proof of it is delivered with it (source0.tex lines 1146--1162, one \texttt{proof} environment of 17 lines). No mathematics is written, repaired or re-derived: the statement and the proof are the article's own lines, in the article's own notation, and no retained line is substituted or re-flowed.  No further source-local context is needed: every cross-reference of every retained line resolves inside the two delivered spans.  Nothing was omitted from any retained span, so the package carries no omission record.  No retained line contains a citation, so the wrapper carries no bibliography and no bibliography entry.  No retained line contains a figure environment, an image-inclusion command or drawing code, so no image is delivered and the package declares no asset dependency.  Editorial formatting only, in addition: an AMS \texttt{amsart} preamble that restates the article's own declarations and macros used by the retained lines, namely the environments \texttt{proposition} on separate counters, together with the standard packages those lines need.  The retained environments print in this document as Proposition 1 in order of appearance, whereas the article prints the same items with the numbering of its own full document.  The complete native source of the article as distributed in the arXiv e-print for this exact version is bundled unchanged as \texttt{sources/200145/source0.tex}.
\begin{proposition}[CV--$\ell_2$ equivalence]
\label{thm:cv-l2}
For any $L\in\mathbb{R}_+^M$ with $\sum_g L_g>0$,
\[
\mathrm{CV}(L)^2 \;=\; M\bigl\|\widehat L\bigr\|_2^2 \;-\; 1.
\]
In particular, minimizing $\mathrm{CV}(L)$ is equivalent to minimizing $\|\widehat L\|_2^2$.
\end{proposition}

\begin{proof}
Let $S:=\sum_{g=1}^M L_g$, so $\mu=S/M$ and $L_g-\mu=S(\widehat L_g-1/M)$.
Then
\[
\sigma^2
= \frac{1}{M}\sum_{g=1}^M S^2\Bigl(\widehat L_g-\frac{1}{M}\Bigr)^2
= \frac{S^2}{M}\sum_{g=1}^M \Bigl(\widehat L_g^2 - \frac{2}{M}\widehat L_g + \frac{1}{M^2}\Bigr).
\]
Since $\sum_g \widehat L_g = 1$, we obtain
$\sum_g (\widehat L_g-\frac{1}{M})^2 = \sum_g \widehat L_g^2 - \frac{1}{M} = \|\widehat L\|_2^2 - \frac{1}{M}$.
\begin{eqnarray*}
\sigma^2
&=& \frac{S^2}{M}\Bigl(\|\widehat L\|_2^2 - \frac{1}{M}\Bigr),\\
\mathrm{CV}(L)^2
&=& M\|\widehat L\|_2^2 - 1.
\end{eqnarray*}
\end{proof}

\end{document}
